%==============================================================================
%  DevPilot AI — Rapport de Projet de Fin d'Études (PFE)
%  Plateforme d'IA d'entreprise pour la gestion des connaissances et le
%  question-réponse fondé sur le RAG (Retrieval-Augmented Generation).
%
%  Compilation :  pdflatex rapport_devpilot_ai.tex   (x2 pour TOC/refs)
%  Aucune dépendance externe (bibliographie en thebibliography).
%  Les zones [INSERER FIGURE] sont des emplacements de captures d'écran.
%==============================================================================
\documentclass[12pt,a4paper,oneside]{report}

%------------------------------------------------------------------ Encodage / langue
\usepackage[utf8]{inputenc}
\usepackage[T1]{fontenc}
\usepackage[french]{babel}
\usepackage{lmodern}
\usepackage{microtype}
\usepackage{csquotes}

%------------------------------------------------------------------ Mise en page
\usepackage[a4paper,margin=2.5cm,headheight=15pt]{geometry}
\usepackage{setspace}
\onehalfspacing
\usepackage{fancyhdr}
\usepackage{parskip}

%------------------------------------------------------------------ Tableaux / figures
\usepackage{graphicx}
\usepackage{booktabs}
\usepackage{array}
\usepackage{tabularx}
\usepackage{longtable}
\usepackage{multirow}
\usepackage{caption}
\usepackage{float}
\usepackage{enumitem}
\usepackage{amsmath}
\usepackage{amssymb}

%------------------------------------------------------------------ Couleurs
\usepackage{xcolor}
\definecolor{brand}{HTML}{6A35F0}
\definecolor{brandDark}{HTML}{4C1FB3}
\definecolor{ink}{HTML}{1A1B25}
\definecolor{slate}{HTML}{475569}
\definecolor{tierClient}{HTML}{EEF2FF}
\definecolor{tierApi}{HTML}{E0F2FE}
\definecolor{tierAsync}{HTML}{FEF3C7}
\definecolor{tierData}{HTML}{DCFCE7}
\definecolor{tierAI}{HTML}{F3E8FF}
\definecolor{codeBg}{HTML}{F6F8FA}
\definecolor{codeKw}{HTML}{8B2FB3}
\definecolor{codeStr}{HTML}{0A7E3D}
\definecolor{codeCom}{HTML}{6B7280}
\definecolor{lineGray}{HTML}{D0D7DE}

%------------------------------------------------------------------ Listings (code)
\usepackage{listings}
\lstdefinestyle{dp}{
  basicstyle=\ttfamily\footnotesize,
  keywordstyle=\color{codeKw}\bfseries,
  stringstyle=\color{codeStr},
  commentstyle=\color{codeCom}\itshape,
  numberstyle=\tiny\color{codeCom},
  backgroundcolor=\color{codeBg},
  frame=single,
  rulecolor=\color{lineGray},
  numbers=left,
  numbersep=8pt,
  tabsize=2,
  breaklines=true,
  breakatwhitespace=true,
  showstringspaces=false,
  columns=fullflexible,
  keepspaces=true,
  captionpos=b,
  xleftmargin=14pt,
  framexleftmargin=14pt,
}
\lstset{style=dp}

%------------------------------------------------------------------ TikZ
\usepackage{tikz}
\usetikzlibrary{positioning,arrows.meta,fit,calc,shapes.geometric,backgrounds,shadows.blur}
\tikzset{
  comp/.style   = {draw=lineGray, line width=0.6pt, rounded corners=2pt, fill=white,
                   align=center, font=\scriptsize, minimum height=0.85cm, inner sep=4pt},
  tier/.style   = {draw=lineGray, line width=0.7pt, rounded corners=3pt,
                   align=center, font=\small\bfseries, inner sep=8pt},
  flow/.style   = {-{Stealth[length=2.4mm]}, line width=0.8pt, color=slate},
  flowb/.style  = {{Stealth[length=2mm]}-{Stealth[length=2mm]}, line width=0.8pt, color=slate},
  store/.style  = {draw=lineGray, line width=0.6pt, cylinder, shape border rotate=90,
                   aspect=0.25, fill=tierData, align=center, font=\scriptsize,
                   minimum height=1cm, minimum width=1.6cm, inner sep=2pt},
  stage/.style  = {draw=lineGray, line width=0.6pt, rounded corners=2pt, fill=tierAI,
                   align=center, font=\scriptsize, minimum height=1.0cm, minimum width=1.9cm},
  life/.style   = {draw=brand, line width=0.7pt, rounded corners=2pt, fill=brand!8,
                   align=center, font=\scriptsize\bfseries, minimum width=2.1cm, minimum height=0.7cm},
}

%------------------------------------------------------------------ Hyperref (dernier)
\usepackage[hidelinks]{hyperref}
\hypersetup{
  colorlinks=true,
  linkcolor=brandDark,
  citecolor=brandDark,
  urlcolor=brandDark,
  pdftitle={DevPilot AI - Rapport de Projet de Fin d'Etudes},
  pdfauthor={DevPilot AI},
}

%------------------------------------------------------------------ En-têtes
\pagestyle{fancy}
\fancyhf{}
\fancyhead[L]{\small\nouppercase{\leftmark}}
\fancyhead[R]{\small\textbf{DevPilot AI}}
\fancyfoot[C]{\thepage}
\renewcommand{\headrulewidth}{0.4pt}
\renewcommand{\footrulewidth}{0pt}

%------------------------------------------------------------------ Raccourcis
\newcommand{\devpilot}{\textsc{DevPilot AI}}
\newcommand{\figplaceholder}[1]{%
  \begin{center}
  \fbox{\parbox[c][3.2cm][c]{0.82\textwidth}{\centering\itshape\color{slate}%
  [INSÉRER FIGURE]\\[2pt]\small #1}}
  \end{center}}
\newcommand{\code}[1]{\texttt{\small #1}}

%==============================================================================
\begin{document}
%==============================================================================

%------------------------------------------------------------------ PAGE DE GARDE
\begin{titlepage}
\centering
\thispagestyle{empty}
{\small\scshape Royaume du Maroc \\ Université Mohammed V de Rabat \\
École Nationale Supérieure d'Informatique et d'Analyse des Systèmes (ENSIAS)\par}
\vspace{0.8cm}
{\color{lineGray}\rule{\textwidth}{0.5pt}}
\vspace{1.2cm}

{\small\itshape Rapport de Projet de Fin d'Études\par}
\vspace{0.4cm}
{\small en vue de l'obtention du diplôme d'Ingénieur d'État\par}
\vspace{0.2cm}
{\small\textit{Filière : Génie Logiciel / Intelligence Artificielle}\par}

\vspace{1.6cm}
{\color{brand}\Huge\bfseries DevPilot AI\par}
\vspace{0.6cm}
{\LARGE\bfseries Conception et réalisation d'une plateforme d'IA\\[4pt]
pour la gestion des connaissances d'ingénierie\\[4pt]
et le question-réponse fondé sur le RAG\par}
\vspace{0.5cm}
{\large\itshape Architecture, pipeline agentique, évaluation et observabilité\par}

\vspace{1.8cm}
{\color{lineGray}\rule{0.6\textwidth}{0.5pt}}
\vspace{0.6cm}

\begin{minipage}{0.46\textwidth}
\begin{flushleft}\large
\textbf{Réalisé par :}\\[2pt]
[Nom et prénom de l'étudiant]
\end{flushleft}
\end{minipage}\hfill
\begin{minipage}{0.46\textwidth}
\begin{flushright}\large
\textbf{Encadré par :}\\[2pt]
[Encadrant pédagogique — ENSIAS]\\
[Encadrant professionnel — Organisme]
\end{flushright}
\end{minipage}

\vspace{1.4cm}
\textbf{Membres du jury :}\\[4pt]
\begin{tabular}{ll}
[Président du jury] & \textit{Président} \\
[Examinateur] & \textit{Examinateur} \\
[Encadrant] & \textit{Encadrant} \\
\end{tabular}

\vfill
{\color{lineGray}\rule{\textwidth}{0.5pt}}
\vspace{0.3cm}
{\small Année universitaire 2025--2026\par}
\end{titlepage}

%------------------------------------------------------------------ Remerciements
\chapter*{Remerciements}
\addcontentsline{toc}{chapter}{Remerciements}
\markboth{Remerciements}{Remerciements}

Au terme de ce projet de fin d'études, je tiens à exprimer ma profonde gratitude
à toutes les personnes qui ont contribué, de près ou de loin, à l'aboutissement
de ce travail.

Je remercie particulièrement mon encadrant pédagogique [\textit{À compléter}] pour
son accompagnement rigoureux, ses conseils méthodologiques et la confiance qu'il
m'a accordée tout au long de ce projet. Mes remerciements s'adressent également à
mon encadrant professionnel [\textit{À compléter}] pour son expertise technique et
sa disponibilité.

Je tiens à remercier l'ensemble du corps professoral de l'ENSIAS pour la qualité
de la formation d'ingénieur dispensée, ainsi que les membres du jury qui ont
accepté d'évaluer ce travail.

Enfin, je dédie ce travail à ma famille, dont le soutien indéfectible a rendu ce
parcours possible.

%------------------------------------------------------------------ Abstract (EN)
\chapter*{Abstract}
\addcontentsline{toc}{chapter}{Abstract}
\markboth{Abstract}{Abstract}

Engineering organizations accumulate large volumes of technical knowledge —
specifications, architecture documents, design notes, internal guidelines and
project reports. Traditional keyword search fails to capture the semantic intent
of a query, while general-purpose generative assistants such as ChatGPT produce
fluent but \emph{ungrounded} answers: they hallucinate, expose no citations, and
offer no visibility into how an answer was produced.

\devpilot{} is an enterprise platform that closes this gap. It implements a
production-grade \emph{Retrieval-Augmented Generation} (RAG) pipeline in which
every answer is grounded in the organization's own documents, supported by
inline citations, scored for faithfulness, and fully traceable end to end. Users
organise their knowledge into isolated \emph{workspaces}, upload documents that
are ingested asynchronously, chunked, embedded and indexed into a vector store,
and then query that knowledge through a streaming conversational interface.

The platform is built as a service-oriented system: a \textsc{FastAPI} backend
exposes a typed REST and streaming API; a \textsc{Celery}/\textsc{Redis}
pipeline performs heavy document processing off the request path;
\textsc{PostgreSQL} stores relational state and the complete execution trace;
and \textsc{Qdrant} provides approximate nearest-neighbour vector search. The
answer-generation stage is an \emph{agentic workflow} of seven cooperating
agents — router, query rewriter, retriever, reranker, generator, evaluator and
corrector — each of which is persisted as a first-class, inspectable run.

A central engineering contribution is \emph{observability}: a RAGOps control
tower monitors pipeline health, and a Trace Inspector reconstructs each query's
full life cycle — agent timeline, retrieved chunks, reranking decisions,
citations, evaluation scores and corrector diff. The result is a system that is
not only useful, but \emph{auditable}, \emph{measurable} and \emph{operable} at
the standard expected of an enterprise AI platform.

\noindent\textbf{Keywords:} Retrieval-Augmented Generation, Large Language
Models, Vector Search, Agentic Pipelines, AI Observability, FastAPI, Qdrant,
Celery, Next.js, Faithfulness Evaluation.

%------------------------------------------------------------------ Résumé (FR)
\chapter*{Résumé}
\addcontentsline{toc}{chapter}{Résumé}
\markboth{Résumé}{Résumé}

Les organisations d'ingénierie accumulent un volume considérable de connaissances
techniques : spécifications, documents d'architecture, notes de conception,
directives internes et rapports de projet. La recherche traditionnelle par
mots-clés échoue à capturer l'intention sémantique d'une requête, tandis que les
assistants génératifs généralistes tels que ChatGPT produisent des réponses
fluides mais \emph{non fondées} : ils hallucinent, n'exposent aucune citation et
n'offrent aucune visibilité sur la manière dont une réponse a été produite.

\devpilot{} est une plateforme d'entreprise qui comble cet écart. Elle met en
{\oe}uvre un pipeline de \emph{Génération Augmentée par Récupération} (RAG) de
qualité production dans lequel chaque réponse est ancrée dans les documents
propres à l'organisation, étayée par des citations en ligne, notée pour sa
fidélité et entièrement traçable de bout en bout. Les utilisateurs organisent
leurs connaissances en \emph{espaces de travail} isolés, téléversent des
documents ingérés de manière asynchrone, découpés, vectorisés et indexés dans
une base vectorielle, puis interrogent cette connaissance à travers une interface
conversationnelle en \emph{streaming}.

La plateforme est conçue comme un système orienté services : un \textit{backend}
\textsc{FastAPI} expose une API REST et \textit{streaming} typée ; un pipeline
\textsc{Celery}/\textsc{Redis} effectue le traitement documentaire lourd hors du
chemin critique des requêtes ; \textsc{PostgreSQL} stocke l'état relationnel et
la trace d'exécution complète ; et \textsc{Qdrant} assure la recherche vectorielle
par plus proches voisins approchés. L'étape de génération repose sur un
\emph{flux agentique} de sept agents coopérants — routeur, réécriture de requête,
récupération, reclassement, génération, évaluation et correction — chacun étant
persisté comme une exécution inspectable de premier ordre.

Une contribution d'ingénierie centrale est l'\emph{observabilité} : une tour de
contrôle RAGOps supervise la santé du pipeline, et un Inspecteur de Traces
reconstruit le cycle de vie complet de chaque requête. Le résultat est un système
non seulement utile, mais \emph{auditable}, \emph{mesurable} et \emph{exploitable}
au niveau attendu d'une plateforme d'IA d'entreprise.

\noindent\textbf{Mots-clés :} Génération Augmentée par Récupération, Grands
Modèles de Langage, Recherche Vectorielle, Pipelines Agentiques, Observabilité de
l'IA, FastAPI, Qdrant, Celery, Next.js, Évaluation de la Fidélité.

%------------------------------------------------------------------ Tables
\tableofcontents
\listoffigures
\listoftables

%==============================================================================
\chapter*{Introduction générale}
\addcontentsline{toc}{chapter}{Introduction générale}
\markboth{Introduction générale}{Introduction générale}
%==============================================================================

\section*{Contexte}
L'intelligence artificielle générative a profondément transformé l'accès à
l'information. Les grands modèles de langage (\emph{Large Language Models}, LLM)
produisent un texte d'une fluidité remarquable et démocratisent l'interaction en
langage naturel. Cependant, dans un contexte d'\textbf{ingénierie d'entreprise},
ces modèles présentent une limite structurelle : ils répondent à partir de leurs
paramètres figés au moment de l'entraînement, sans connaissance du patrimoine
documentaire propre à l'organisation, et sans aucune garantie de véracité. Pour
une équipe d'ingénierie, une réponse plausible mais fausse — une
\emph{hallucination} — n'est pas une simple imperfection : elle constitue un
risque opérationnel.

\section*{Motivation}
Les équipes techniques produisent et consomment en permanence une documentation
hétérogène et fragmentée. Retrouver l'information pertinente exige aujourd'hui un
effort cognitif disproportionné, et la recherche par mots-clés des outils
classiques ne comprend pas l'intention de la question. La promesse de la
\emph{Génération Augmentée par Récupération} (RAG) est précisément de réconcilier
la puissance générative des LLM avec la connaissance vérifiable de l'organisation.
Mais transformer cette promesse en \emph{système de production} suppose de
résoudre des problèmes d'ingénierie réels : ingestion asynchrone, recherche
vectorielle, ancrage et citation des sources, mesure de la qualité, et surtout
\emph{observabilité} de bout en bout.

\section*{Problématique}
Comment concevoir et réaliser une plateforme d'IA capable de \textbf{centraliser}
les connaissances d'une équipe d'ingénierie, de \textbf{récupérer} l'information
pertinente, de \textbf{générer} des réponses fondées et \textbf{citées}, de
\textbf{mesurer} la qualité de ces réponses, et d'offrir une \textbf{traçabilité}
complète permettant d'auditer et d'exploiter le système comme tout service
critique d'entreprise ?

\section*{Objectifs}
Le présent projet vise à concevoir et implémenter \devpilot{}, une plateforme RAG
multi-tenant offrant : la gestion d'espaces de travail isolés ; l'ingestion
documentaire asynchrone (découpage, vectorisation, indexation Qdrant) ; la
récupération sémantique et hybride avec reclassement ; la génération agentique de
réponses ancrées et citées ; l'évaluation automatique de la fidélité et la
détection d'hallucinations ; le \emph{streaming} temps réel ; et un sous-système
d'observabilité (RAGOps, explorateur et inspecteur de traces).

\section*{Méthodologie}
La démarche adoptée est itérative et orientée production. Elle articule une phase
d'analyse des besoins (chapitre~1), une phase de conception fondée sur la
modélisation UML et la conception de données (chapitre~2), une phase
d'\emph{architecture système} qui constitue le c{\oe}ur du travail
(chapitre~3), puis une phase de réalisation et de validation (chapitres~4 et~5).
Chaque décision technique est justifiée par une analyse de compromis explicite.

\section*{Organisation du rapport}
Le chapitre~1 analyse le contexte et les besoins. Le chapitre~2 présente la
conception (UML, modèle de données, traçabilité). Le chapitre~3, le plus étendu,
détaille l'architecture système, les pipelines et les décisions d'ingénierie. Le
chapitre~4 décrit la réalisation technique. Le chapitre~5 traite des tests et de
la validation. Le rapport se clôt par les limites, les perspectives et une
conclusion générale.

%==============================================================================
\chapter{Contexte général et analyse des besoins}
\label{chap:contexte}
%==============================================================================

\section{Introduction}
Ce chapitre établit les fondations du projet. Il caractérise le problème métier,
positionne \devpilot{} face aux solutions existantes, puis formalise les exigences
fonctionnelles et non fonctionnelles ainsi que les contraintes du projet. Cette
analyse rigoureuse conditionne l'ensemble des choix d'architecture présentés au
chapitre~3.

\section{Contexte et motivation}
La connaissance d'ingénierie possède trois propriétés qui la rendent difficile à
exploiter. Elle est d'abord \textbf{volumineuse et croissante} : chaque projet
génère des centaines de pages de spécifications, de comptes rendus et de
documentation technique. Elle est ensuite \textbf{fragmentée} : répartie entre
des outils, des formats et des silos hétérogènes. Elle est enfin
\textbf{contextuelle} : une même question n'a pas la même réponse selon le projet,
l'équipe ou la version. Un assistant générique, entraîné sur le web public,
ignore par construction ce contexte propriétaire.

\section{Énoncé du problème}
Le problème se décompose en deux familles de limites complémentaires.

\paragraph{Limites de la recherche traditionnelle.} Les moteurs de recherche
internes reposent majoritairement sur l'appariement lexical (mots-clés). Ils
échouent dès que le vocabulaire de la question diffère de celui du document, ne
comprennent pas la sémantique, et renvoient des listes de documents que
l'utilisateur doit encore lire et synthétiser lui-même.

\paragraph{Limites des assistants génératifs.} Les LLM généralistes produisent une
synthèse directe, mais souffrent d'\textbf{hallucinations}, d'une \textbf{absence
de traçabilité}, d'une \textbf{absence de citations}, d'une \textbf{absence de
visibilité} sur la qualité de récupération, et d'une \textbf{absence
d'observabilité}. Ils ne peuvent pas, en l'état, être tenus pour responsables de
leurs réponses.

\section{Objectifs du projet}
\devpilot{} se fixe pour objectif de construire un système qui réconcilie ces deux
mondes. Les objectifs opérationnels sont synthétisés dans le tableau~\ref{tab:objectifs}.

\begin{table}[H]
\centering
\caption{Objectifs opérationnels et capacités correspondantes}
\label{tab:objectifs}
\small
\begin{tabularx}{\textwidth}{@{}>{\bfseries}l X@{}}
\toprule
Objectif & Capacité réalisée \\
\midrule
Centraliser & Espaces de travail isolés, gestion documentaire, ingestion asynchrone \\
Récupérer & Recherche sémantique et hybride, reclassement par pertinence \\
Générer & Réponses ancrées produites par un pipeline agentique \\
Citer & Génération de citations reliées aux passages sources \\
Mesurer & Évaluation de fidélité, pertinence, détection d'hallucinations \\
Diffuser & \emph{Streaming} temps réel des réponses (SSE) \\
Superviser & RAGOps : santé des pipelines, file de traitement, alertes \\
Tracer & Explorateur et inspecteur de traces de bout en bout \\
\bottomrule
\end{tabularx}
\end{table}

\section{Analyse des solutions existantes}
Aucune solution unique du marché ne couvre simultanément les axes
\emph{ancrage documentaire}, \emph{traçabilité} et \emph{observabilité} dans un
contexte auto-hébergeable. Le tableau~\ref{tab:comparaison} positionne \devpilot{}
par rapport à cinq familles de solutions représentatives.

\begin{table}[H]
\centering
\caption{Analyse comparative des solutions existantes}
\label{tab:comparaison}
\scriptsize
\renewcommand{\arraystretch}{1.25}
\begin{tabularx}{\textwidth}{@{}>{\raggedright\arraybackslash}p{3.0cm} *{6}{>{\centering\arraybackslash}X}@{}}
\toprule
\textbf{Capacité} & \textbf{ChatGPT} & \textbf{Notion AI} & \textbf{LangSmith} & \textbf{Datadog} & \textbf{Recherche interne} & \textbf{\devpilot{}} \\
\midrule
Réponses fondées sur vos documents      & Limité & Oui     & N/A    & Non & Partiel & \textbf{Oui} \\
Citations des sources                    & Non    & Partiel & N/A    & Non & Non     & \textbf{Oui} \\
Traçabilité complète des requêtes        & Non    & Non     & Oui    & Oui & Non     & \textbf{Oui} \\
Isolation par espace de travail          & Non    & Partiel & N/A    & N/A & Partiel & \textbf{Oui} \\
Score d'hallucination                    & Non    & Non     & Partiel& Non & Non     & \textbf{Oui} \\
Métriques de récupération                & Non    & Non     & Oui    & Non & Non     & \textbf{Oui} \\
Chronologie d'exécution des agents       & Non    & Non     & Oui    & Oui & Non     & \textbf{Oui} \\
Évaluation intégrée de la qualité        & Non    & Non     & Oui    & Non & Non     & \textbf{Oui} \\
Déploiement sur site (\emph{on-premise}) & Non    & Non     & Partiel& Oui & Oui     & \textbf{Oui} \\
\bottomrule
\end{tabularx}
\end{table}

\noindent\textbf{Lecture.} ChatGPT et Notion AI excellent sur l'expérience
conversationnelle mais n'offrent ni traçabilité ni observabilité opérationnelle.
LangSmith et Datadog sont des outils d'\emph{observabilité} mais ne constituent
pas une plateforme RAG complète. Les moteurs internes restent lexicaux.
\devpilot{} se distingue en \textbf{intégrant} l'application RAG \emph{et} son
observabilité dans un système unique, auto-hébergeable.

\section{Exigences fonctionnelles}
Les exigences fonctionnelles (EF) sont organisées par domaine et résumées dans le
tableau~\ref{tab:ef}.

\begin{table}[H]
\centering
\caption{Exigences fonctionnelles principales}
\label{tab:ef}
\small
\begin{tabularx}{\textwidth}{@{}l l X@{}}
\toprule
\textbf{Réf.} & \textbf{Domaine} & \textbf{Exigence} \\
\midrule
EF-1  & Authentification & Inscription, connexion et sessions sécurisées par jetons JWT \\
EF-2  & Autorisation     & Contrôle d'accès par rôles (utilisateur, admin, super-admin) \\
EF-3  & Espaces          & Création d'espaces de travail isolés et gestion des membres \\
EF-4  & Documents        & Téléversement, listing, suivi de statut et suppression de documents \\
EF-5  & Ingestion        & Traitement asynchrone : extraction, découpage, vectorisation, indexation \\
EF-6  & Récupération     & Recherche sémantique et hybride sur la base vectorielle \\
EF-7  & Reclassement     & Reclassement des candidats par pertinence (\emph{reranking}) \\
EF-8  & Génération       & Production de réponses ancrées par un pipeline agentique \\
EF-9  & Citations        & Génération de citations reliées aux passages utilisés \\
EF-10 & Conversation     & Historique de conversations et de messages persistants \\
EF-11 & Streaming        & Diffusion incrémentale des réponses (jeton par jeton) \\
EF-12 & Évaluation       & Mesure de la fidélité, de la pertinence et du score d'hallucination \\
EF-13 & Correction       & Réparation automatique des réponses non fidèles (agent correcteur) \\
EF-14 & Administration   & Tableaux de bord d'administration (utilisateurs, espaces, documents) \\
EF-15 & RAGOps           & Supervision de la santé des pipelines et de la file de traitement \\
EF-16 & Traçabilité      & Explorateur et inspecteur de traces de requêtes RAG \\
\bottomrule
\end{tabularx}
\end{table}

\section{Exigences non fonctionnelles}
Les exigences non fonctionnelles (ENF) définissent les qualités attendues du
système ; elles sont déterminantes pour une plateforme d'entreprise et structurent
le chapitre~3.

\begin{description}[leftmargin=2.2cm,style=nextline]
\item[Sécurité] Authentification JWT, autorisation par rôles, isolation stricte
des données par espace de travail, journalisation d'audit des actions
d'administration, et masquage systématique des secrets dans les vues de débogage.
\item[Performance] API asynchrone non bloquante, déport du traitement lourd hors
du chemin critique via Celery, recherche vectorielle par plus proches voisins
approchés (HNSW), et \emph{time-to-first-token} réduit grâce au \emph{streaming}.
\item[Scalabilité] \textit{Backend} sans état, horizontalement extensible ;
\emph{workers} Celery ajoutables à la demande ; base vectorielle Qdrant
partitionnable.
\item[Disponibilité] Conteneurisation de chaque service, dépendances de santé
explicites, et politiques de relance des tâches asynchrones.
\item[Maintenabilité] Architecture en couches, typage statique de bout en bout
(TypeScript côté \textit{frontend}, Pydantic/SQLAlchemy côté \textit{backend}),
migrations de schéma versionnées (Alembic).
\item[Observabilité] Traçabilité complète des requêtes, métriques Prometheus,
tableaux de bord Grafana, RAGOps et inspecteur de traces.
\item[Fiabilité] Mécanismes de repli de modèles (\emph{fallback}), agent
correcteur, et relance automatique des ingestions échouées.
\item[Utilisabilité] Interface conversationnelle réactive, rendu Markdown des
réponses, cartes de citation interactives, thème clair/sombre.
\end{description}

\section{Contraintes}
\begin{description}[leftmargin=2.4cm,style=nextline]
\item[Techniques] Reproductibilité par conteneurs ; possibilité d'exécution
hors-ligne via des modèles locaux (Ollama) ; budget mémoire et calcul maîtrisé.
\item[Académiques] Démonstration d'une démarche d'ingénierie rigoureuse et d'un
raisonnement architectural, au-delà d'un simple prototype.
\item[Temps] Réalisation dans le cadre temporel d'un projet de fin d'études.
\item[Ressources] Environnement de développement et d'exécution mono-machine,
sans infrastructure cloud dédiée.
\end{description}

\section{Cas d'utilisation}
Les acteurs principaux sont l'\textbf{utilisateur} (membre d'une équipe
d'ingénierie), l'\textbf{administrateur} (gestion de la plateforme) et le
\textbf{super-administrateur} (observabilité approfondie). Les cas d'utilisation
clés incluent : gérer un espace de travail, téléverser et suivre un document,
poser une question et recevoir une réponse citée en \emph{streaming}, consulter
l'historique, superviser les pipelines (RAGOps) et inspecter une trace. Leur
formalisation UML est présentée au chapitre~2.

\section{Conclusion du chapitre}
Ce chapitre a établi que le besoin n'est pas un \emph{chatbot} de plus, mais une
plateforme d'IA \emph{fondée}, \emph{mesurable} et \emph{observable}. Les exigences
non fonctionnelles — sécurité, scalabilité, observabilité — sont au moins aussi
structurantes que les exigences fonctionnelles. Le chapitre suivant traduit ces
besoins en modèles de conception.

%==============================================================================
\chapter{Analyse et conception}
\label{chap:conception}
%==============================================================================

\section{Introduction}
Ce chapitre formalise la conception du système : modélisation des cas
d'utilisation, modèle du domaine, modèle de données relationnel, scénarios
d'interaction (diagrammes de séquence) et — point distinctif du projet — le
modèle de \emph{traçabilité} qui relie chaque message à son exécution complète.

\section{Diagramme de cas d'utilisation global}
La figure~\ref{fig:usecase} présente la vue d'ensemble des cas d'utilisation et
de leurs acteurs. Le super-administrateur hérite des capacités de
l'administrateur, qui hérite à son tour de celles de l'utilisateur.

\begin{figure}[H]
\centering
\resizebox{0.95\textwidth}{!}{%
\begin{tikzpicture}[every node/.style={font=\scriptsize}]
  % Acteurs
  \node (u)  at (0,0)    {\begin{tabular}{c}\textbf{Utilisateur}\end{tabular}};
  \node (a)  at (0,-3)   {\begin{tabular}{c}\textbf{Administrateur}\end{tabular}};
  \node (s)  at (0,-6)   {\begin{tabular}{c}\textbf{Super-admin}\end{tabular}};
  \foreach \n in {u,a,s}{ \node[draw,circle,minimum size=4mm] at ([yshift=7mm]\n.north) {}; }
  % Système (boite)
  \node[draw,rounded corners,minimum width=8.5cm,minimum height=8cm,label=above:\textbf{Plateforme DevPilot AI}] (sys) at (6.5,-2.6) {};
  % Cas d'utilisation
  \node[draw,ellipse,minimum width=3cm] (uc1) at (5.2,0.6)  {Gérer un espace};
  \node[draw,ellipse,minimum width=3cm] (uc2) at (8.0,-0.4) {Téléverser un document};
  \node[draw,ellipse,minimum width=3cm] (uc3) at (5.0,-1.4) {Poser une question (RAG)};
  \node[draw,ellipse,minimum width=3cm] (uc4) at (8.2,-2.4) {Consulter l'historique};
  \node[draw,ellipse,minimum width=3cm] (uc5) at (5.0,-3.4) {Administrer les utilisateurs};
  \node[draw,ellipse,minimum width=3cm] (uc6) at (8.0,-4.2) {Superviser (RAGOps)};
  \node[draw,ellipse,minimum width=3cm] (uc7) at (5.2,-5.2) {Inspecter une trace};
  % Associations
  \draw (u) -- (uc1); \draw (u) -- (uc2); \draw (u) -- (uc3); \draw (u) -- (uc4);
  \draw (a) -- (uc5); \draw (a) -- (uc6);
  \draw (s) -- (uc7);
  % Heritage acteurs
  \draw[-{Stealth},dashed] (a) -- node[right,font=\tiny]{$\langle$hérite$\rangle$} (u);
  \draw[-{Stealth},dashed] (s) -- node[right,font=\tiny]{$\langle$hérite$\rangle$} (a);
\end{tikzpicture}}
\caption{Diagramme de cas d'utilisation global}
\label{fig:usecase}
\end{figure}

\section{Cas d'utilisation détaillés}
Le tableau~\ref{tab:ucdetail} détaille le cas d'utilisation pivot
« Poser une question (RAG) », dont le scénario nominal active l'ensemble du
pipeline agentique.

\begin{table}[H]
\centering
\caption{Cas d'utilisation détaillé : « Poser une question (RAG) »}
\label{tab:ucdetail}
\small
\begin{tabularx}{\textwidth}{@{}>{\bfseries}l X@{}}
\toprule
Acteur & Utilisateur authentifié, membre d'un espace de travail \\
Pré-condition & L'espace contient au moins un document indexé \\
Déclencheur & L'utilisateur soumet une question dans l'interface de chat \\
\midrule
\multirow{1}{*}{Scénario nominal} & 1.~Le routeur valide et classe la requête.\newline
2.~La requête est réécrite pour améliorer la récupération.\newline
3.~Les passages candidats sont récupérés depuis Qdrant.\newline
4.~Les candidats sont reclassés par pertinence.\newline
5.~Le générateur produit une réponse ancrée et citée.\newline
6.~L'évaluateur note la fidélité et l'hallucination.\newline
7.~Si nécessaire, le correcteur répare la réponse.\newline
8.~La réponse est diffusée en \emph{streaming} avec ses citations. \\
\midrule
Post-condition & Le message, ses agents, ses passages, son évaluation sont persistés \\
Exceptions & Aucun passage pertinent $\rightarrow$ réponse de repli explicite ;
échec LLM $\rightarrow$ modèle de \emph{fallback} \\
\bottomrule
\end{tabularx}
\end{table}

\section{Modèle du domaine}
Le modèle du domaine articule sept concepts métier centraux : l'\textbf{Utilisateur}
appartient à des \textbf{Espaces de travail} via des \textbf{Adhésions} ; un espace
contient des \textbf{Documents} découpés en \textbf{Chunks} ; un utilisateur mène
des \textbf{Conversations} composées de \textbf{Messages} ; chaque message
assistant agrège une \textbf{Trace} (agents, passages récupérés, évaluation). Ce
dernier point — élever la trace au rang de concept du domaine — est un choix de
conception déterminant pour l'observabilité.

\section{Modèle de données}
Le schéma relationnel comprend onze tables, implémentées avec SQLAlchemy et
versionnées par Alembic. La figure~\ref{fig:erd} en présente une vue entité-association
simplifiée, et le tableau~\ref{tab:tables} en décrit le rôle.

\begin{figure}[H]
\centering
\resizebox{\textwidth}{!}{%
\begin{tikzpicture}[
  ent/.style={draw=lineGray,rounded corners=2pt,fill=tierApi,font=\scriptsize\bfseries,
              minimum width=2.3cm,minimum height=0.7cm},
  rel/.style={-{Stealth[length=2mm]},line width=0.7pt,color=slate},
  lbl/.style={font=\tiny\itshape,color=slate,fill=white,inner sep=1pt}]
  \node[ent] (users)   at (0,0)    {users};
  \node[ent] (wm)      at (3.4,0)  {workspace\_members};
  \node[ent] (ws)      at (7,0)    {workspaces};
  \node[ent] (docs)    at (10.4,0) {documents};
  \node[ent] (chunks)  at (10.4,-2){chunks};
  \node[ent] (conv)    at (0,-2)   {conversations};
  \node[ent] (msg)     at (3.4,-2) {messages};
  \node[ent] (ar)      at (3.4,-4) {agent\_runs};
  \node[ent] (rc)      at (0,-4)   {retrieved\_chunks};
  \node[ent] (ev)      at (6.8,-4) {evaluations};
  \node[ent] (lu)      at (10.4,-4){llm\_usage};
  % relations
  \draw[rel] (users) -- node[lbl]{1..*} (wm);
  \draw[rel] (ws)    -- node[lbl]{1..*} (wm);
  \draw[rel] (ws)    -- node[lbl]{1..*} (docs);
  \draw[rel] (docs)  -- node[lbl]{1..*} (chunks);
  \draw[rel] (users) -- node[lbl]{1..*} (conv);
  \draw[rel] (ws)    -- node[lbl]{0..*} (conv);
  \draw[rel] (conv)  -- node[lbl]{1..*} (msg);
  \draw[rel] (msg)   -- node[lbl]{1..*} (ar);
  \draw[rel] (msg)   -- node[lbl]{1..*} (rc);
  \draw[rel] (msg)   -- node[lbl]{1..1} (ev);
  \draw[rel] (msg)   -- node[lbl]{0..*} (lu);
\end{tikzpicture}}
\caption{Modèle de données — vue entité-association simplifiée (11 tables)}
\label{fig:erd}
\end{figure}

\begin{table}[H]
\centering
\caption{Description des tables du modèle de données}
\label{tab:tables}
\small
\begin{tabularx}{\textwidth}{@{}>{\ttfamily}l X@{}}
\toprule
\normalfont\textbf{Table} & \textbf{Rôle} \\
\midrule
users              & Comptes, identifiants, rôle (utilisateur/admin/super-admin) \\
workspaces         & Espaces de travail isolant les connaissances \\
workspace\_members & Adhésions reliant utilisateurs et espaces (table d'association) \\
documents          & Documents téléversés et leur statut d'ingestion \\
chunks             & Fragments de texte avec référence vers le vecteur Qdrant \\
conversations      & Fils de discussion d'un utilisateur dans un espace \\
messages           & Messages utilisateur et assistant d'une conversation \\
agent\_runs        & Exécution de chaque agent du pipeline (statut, latence, E/S) \\
retrieved\_chunks  & Passages récupérés et associés à un message assistant \\
evaluations        & Scores de fidélité, pertinence, hallucination par message \\
llm\_usage         & Comptabilité d'usage des modèles (jetons, modèle, coût) \\
\bottomrule
\end{tabularx}
\end{table}

\section{Diagrammes de séquence}
Quatre scénarios d'interaction structurent le système. Faute de place, ils sont
présentés sous forme de diagrammes de séquence simplifiés (lignes de vie et
messages). Les retours synchrones sont en pointillés.

\subsection{Ingestion d'un document}
\begin{figure}[H]
\centering
\resizebox{0.96\textwidth}{!}{%
\begin{tikzpicture}[font=\scriptsize,>=Stealth]
  \def\h{-6.2}
  \foreach \x/\name in {0/Client,3/API,6/Stockage,9/Celery,12/Qdrant}{
    \node[life] at (\x,0) {\name};
    \draw[dashed,lineGray] (\x,-0.4) -- (\x,\h);}
  \draw[->] (0,-1)  -- (3,-1)  node[midway,above]{POST /documents (fichier)};
  \draw[->] (3,-1.7)-- (6,-1.7)node[midway,above]{persister fichier + métadonnées};
  \draw[->] (3,-2.4)-- (9,-2.4)node[midway,above]{enqueue process\_document\_task};
  \draw[->,dashed] (3,-3.0)-- (0,-3.0) node[midway,above]{202 Accepted (statut=pending)};
  \draw[->] (9,-3.8)-- (6,-3.8)node[midway,above]{lire + extraire le texte};
  \draw[->] (9,-4.5)-- (9.0,-4.9) ;
  \node[font=\tiny,align=center] at (9,-4.4) {découper + vectoriser};
  \draw[->] (9,-5.3)-- (12,-5.3)node[midway,above]{upsert vecteurs (768-d)};
  \draw[->] (9,-6.0)-- (6,-6.0)node[midway,above]{statut=indexed};
\end{tikzpicture}}
\caption{Diagramme de séquence — ingestion asynchrone d'un document}
\label{fig:seq-ingest}
\end{figure}

\subsection{Requête RAG}
\begin{figure}[H]
\centering
\resizebox{0.96\textwidth}{!}{%
\begin{tikzpicture}[font=\scriptsize,>=Stealth]
  \def\h{-7.4}
  \foreach \x/\name in {0/Client,3/API,6/Workflow,9/Qdrant,12/LLM}{
    \node[life] at (\x,0) {\name};
    \draw[dashed,lineGray] (\x,-0.4) -- (\x,\h);}
  \draw[->] (0,-1)  -- (3,-1)  node[midway,above]{POST /chat (question)};
  \draw[->] (3,-1.6)-- (6,-1.6)node[midway,above]{exécuter le pipeline agentique};
  \draw[->] (6,-2.3)-- (6.0,-2.6); \node[font=\tiny] at (6,-2.15){router + réécriture};
  \draw[->] (6,-3.1)-- (9,-3.1)node[midway,above]{récupérer candidats (k=15)};
  \draw[->,dashed] (9,-3.7)-- (6,-3.7)node[midway,above]{passages + scores};
  \draw[->] (6,-4.3)-- (6.0,-4.6); \node[font=\tiny] at (6,-4.15){reclasser (top 5)};
  \draw[->] (6,-5.1)-- (12,-5.1)node[midway,above]{générer réponse ancrée};
  \draw[->,dashed] (12,-5.7)-- (6,-5.7)node[midway,above]{réponse + citations};
  \draw[->] (6,-6.3)-- (6.0,-6.6); \node[font=\tiny] at (6,-6.15){évaluer + corriger};
  \draw[->,dashed] (6,-7.1)-- (0,-7.1)node[midway,above]{réponse finale + trace persistée};
\end{tikzpicture}}
\caption{Diagramme de séquence — requête RAG (pipeline agentique complet)}
\label{fig:seq-rag}
\end{figure}

\subsection{Requête en streaming}
\begin{figure}[H]
\centering
\resizebox{0.92\textwidth}{!}{%
\begin{tikzpicture}[font=\scriptsize,>=Stealth]
  \def\h{-5.4}
  \foreach \x/\name in {0/Client,4/API SSE,8/Workflow}{
    \node[life] at (\x,0) {\name};
    \draw[dashed,lineGray] (\x,-0.4) -- (\x,\h);}
  \draw[->] (0,-1)  -- (4,-1)  node[midway,above]{POST /chat/stream};
  \draw[->] (4,-1.6)-- (8,-1.6)node[midway,above]{démarrer le pipeline};
  \draw[->,dashed] (8,-2.2)-- (4,-2.2)node[midway,above]{event: trace / citations};
  \draw[->,dashed] (4,-2.8)-- (0,-2.8)node[midway,above]{flux SSE: token, token, ...};
  \draw[->,dashed] (8,-3.6)-- (4,-3.6)node[midway,above]{event: evaluation};
  \draw[->,dashed] (4,-4.2)-- (0,-4.2)node[midway,above]{event: message (final)};
  \draw[->,dashed] (4,-4.9)-- (0,-4.9)node[midway,above]{event: done};
\end{tikzpicture}}
\caption{Diagramme de séquence — diffusion en \emph{streaming} (Server-Sent Events)}
\label{fig:seq-stream}
\end{figure}

\subsection{Supervision administrateur}
\begin{figure}[H]
\centering
\resizebox{0.92\textwidth}{!}{%
\begin{tikzpicture}[font=\scriptsize,>=Stealth]
  \def\h{-5.0}
  \foreach \x/\name in {0/Admin,4/API,8/PostgreSQL,11.5/Qdrant}{
    \node[life] at (\x,0) {\name};
    \draw[dashed,lineGray] (\x,-0.4) -- (\x,\h);}
  \draw[->] (0,-1)  -- (4,-1)  node[midway,above]{GET /admin/ragops/...};
  \draw[->] (4,-1.7)-- (8,-1.7)node[midway,above]{agréger états documents/traces};
  \draw[->] (4,-2.4)-- (11.5,-2.4)node[midway,above]{santé de la collection};
  \draw[->,dashed] (8,-3.1)-- (4,-3.1)node[midway,above]{métriques + file};
  \draw[->,dashed] (4,-3.8)-- (0,-3.8)node[midway,above]{tableau de bord RAGOps};
  \draw[->] (0,-4.5)-- (4,-4.5)node[midway,above]{GET /admin/rag-traces/{id}};
\end{tikzpicture}}
\caption{Diagramme de séquence — supervision et inspection (RAGOps)}
\label{fig:seq-admin}
\end{figure}

\section{Modèle de traçabilité}
La traçabilité est modélisée comme une chaîne dirigée centrée sur le message
assistant : \[\texttt{message} \rightarrow \texttt{agent\_runs} \rightarrow
\texttt{retrieved\_chunks} \rightarrow \texttt{evaluations}.\] Chaque message
assistant agrège l'ensemble des exécutions d'agents qui l'ont produit, les
passages effectivement récupérés, et son évaluation. Cette modélisation rend
chaque réponse \emph{reconstructible a posteriori} sans dépendance à des journaux
externes, et fonde directement l'inspecteur de traces du chapitre~4.

\section{Décisions de conception}
Trois décisions de conception méritent d'être soulignées. Premièrement,
\textbf{la trace est un citoyen de premier ordre} du modèle de données, et non un
artefact de journalisation : ce choix est la condition de l'observabilité.
Deuxièmement, \textbf{l'ingestion est asynchrone} par construction, séparant le
chemin critique de l'API du traitement lourd. Troisièmement,
\textbf{la génération est agentique}, chaque étape étant isolée, persistée et
indépendamment mesurable.

\section{Conclusion du chapitre}
La conception traduit les exigences en modèles exploitables : cas d'utilisation,
modèle de données à onze tables, scénarios d'interaction et chaîne de traçabilité.
Ces modèles préparent l'architecture système détaillée au chapitre suivant.

%==============================================================================
\chapter{System Design et architecture}
\label{chap:archi}
%==============================================================================

\section{Introduction et vision architecturale}
Ce chapitre constitue le c{\oe}ur du travail d'ingénierie. La vision
architecturale de \devpilot{} repose sur quatre principes directeurs :
(i)~\textbf{séparation des responsabilités} en couches et services autonomes ;
(ii)~\textbf{asynchronisme} pour découpler le chemin critique des requêtes du
traitement lourd ; (iii)~\textbf{traçabilité native}, chaque étape produisant une
trace persistée et inspectable ; et (iv)~\textbf{neutralité des fournisseurs}
d'IA, le système s'abstrayant des modèles concrets (Ollama, Gemini, OpenAI). Ces
principes guident l'ensemble des décisions présentées ci-après.

\section{Architecture système globale}
\devpilot{} est un système orienté services déployé par conteneurs. La
figure~\ref{fig:archi-global} présente la topologie complète : huit services
coopérants, organisés en quatre niveaux logiques (présentation, application,
traitement asynchrone, données et IA).

\begin{figure}[H]
\centering
\resizebox{\textwidth}{!}{%
\begin{tikzpicture}[node distance=6mm]
  % Niveau présentation
  \node[comp,fill=tierClient,minimum width=4cm] (next) {Frontend Next.js\\(React 19, Tailwind, React Query)};
  % Niveau application
  \node[comp,fill=tierApi,minimum width=4cm,below=14mm of next] (api) {Backend FastAPI\\(REST + SSE, JWT, services)};
  % stores
  \node[store,below=16mm of api,xshift=-5.2cm] (pg)   {PostgreSQL};
  \node[store,below=16mm of api,xshift=-1.7cm] (redis){Redis};
  \node[store,below=16mm of api,xshift=1.7cm]  (qd)   {Qdrant};
  % async
  \node[comp,fill=tierAsync,minimum width=3cm,below=16mm of api,xshift=5.3cm] (cel) {Celery\\worker};
  % IA
  \node[comp,fill=tierAI,minimum width=2.6cm,below=14mm of qd,xshift=-1.6cm] (ollama){Ollama\\(embeddings)};
  \node[comp,fill=tierAI,minimum width=2.6cm,below=14mm of qd,xshift=1.6cm]  (gemini){Gemini\\(génération)};
  % observabilité
  \node[comp,minimum width=2.6cm,right=10mm of api] (prom){Prometheus};
  \node[comp,minimum width=2.6cm,right=8mm of prom] (graf){Grafana};
  % liens
  \draw[flow] (next) -- node[right,font=\tiny]{HTTPS / SSE} (api);
  \draw[flow] (api) -- (pg);
  \draw[flow] (api) -- (redis);
  \draw[flow] (api) -- (qd);
  \draw[flow] (api) -- (cel);
  \draw[flow] (redis) -- node[above,font=\tiny]{broker} (cel);
  \draw[flow] (cel) |- (pg);
  \draw[flow] (cel) -- (qd);
  \draw[flow] (api) -- (ollama);
  \draw[flow] (api) -- (gemini);
  \draw[flow] (cel) -- (ollama);
  \draw[flow] (api) -- (prom);
  \draw[flow] (prom) -- (graf);
\end{tikzpicture}}
\caption{Architecture système globale — huit services conteneurisés}
\label{fig:archi-global}
\end{figure}

\section{Architecture en couches}
Au-delà de la topologie physique, le système s'organise en couches logiques aux
responsabilités strictes (figure~\ref{fig:archi-layers}). Cette stratification
garantit que les dépendances ne circulent que du haut vers le bas, condition de
la maintenabilité.

\begin{figure}[H]
\centering
\resizebox{0.92\textwidth}{!}{%
\begin{tikzpicture}[node distance=4mm]
  \node[tier,fill=tierClient,minimum width=12cm] (l1)
    {Couche présentation \\ \normalfont\scriptsize Pages App Router, composants UI, chat \emph{streaming}, tableaux de bord d'administration};
  \node[tier,fill=tierApi,minimum width=12cm,below=of l1] (l2)
    {Couche API / Application \\ \normalfont\scriptsize Routes FastAPI, authentification JWT, services métier, schémas Pydantic};
  \node[tier,fill=tierAI,minimum width=12cm,below=of l2] (l3)
    {Couche domaine IA \\ \normalfont\scriptsize Pipeline agentique, récupération, reclassement, génération, évaluation};
  \node[tier,fill=tierAsync,minimum width=12cm,below=of l3] (l4)
    {Couche traitement asynchrone \\ \normalfont\scriptsize Tâches Celery, file Redis, ingestion documentaire};
  \node[tier,fill=tierData,minimum width=12cm,below=of l4] (l5)
    {Couche données \\ \normalfont\scriptsize PostgreSQL (état + traces), Qdrant (vecteurs), stockage de fichiers};
  \foreach \a/\b in {l1/l2,l2/l3,l3/l4,l4/l5}{ \draw[flowb] (\a) -- (\b); }
\end{tikzpicture}}
\caption{Architecture logique en couches}
\label{fig:archi-layers}
\end{figure}

\section{Pipeline de traitement documentaire}
L'ingestion documentaire est entièrement asynchrone. À la réception d'un fichier,
l'API persiste le document avec le statut \code{pending} et délègue le traitement
à Celery via la tâche \code{process\_document\_task} (relancée jusqu'à trois fois
en cas d'échec). La figure~\ref{fig:pipe-doc} décrit les étapes.

\begin{figure}[H]
\centering
\resizebox{\textwidth}{!}{%
\begin{tikzpicture}[node distance=5mm,>=Stealth]
  \node[stage,fill=tierAsync] (up)  {Téléversement\\statut=pending};
  \node[stage,fill=tierAsync,right=of up] (ex)  {Extraction\\du texte};
  \node[stage,fill=tierAsync,right=of ex] (ch)  {Découpage\\(chunks)};
  \node[stage,fill=tierAI,right=of ch]    (em)  {Vectorisation\\(768-d)};
  \node[stage,fill=tierData,right=of em]  (ix)  {Indexation\\Qdrant};
  \node[stage,fill=tierData,right=of ix]  (ok)  {statut=indexed};
  \foreach \a/\b in {up/ex,ex/ch,ch/em,em/ix,ix/ok}{ \draw[flow] (\a) -- (\b); }
  \node[font=\tiny\itshape,below=8mm of ch,color=slate] {Exécuté par le \emph{worker} Celery, hors du chemin critique de l'API};
\end{tikzpicture}}
\caption{Pipeline de traitement documentaire (ingestion asynchrone)}
\label{fig:pipe-doc}
\end{figure}

\section{Pipeline d'\emph{embedding} et architecture de recherche vectorielle}
Les fragments de texte sont vectorisés par le modèle \code{nomic-embed-text}
(Ollama), produisant des vecteurs de dimension~768, traités par lots de taille
configurable. Chaque vecteur est inséré (\emph{upsert}) dans la collection Qdrant,
tandis que PostgreSQL conserve le fragment et une référence \code{vector\_id} vers
le point Qdrant correspondant. Ce \textbf{couplage dual} — texte canonique en
relationnel, vecteur en base vectorielle — sépare les responsabilités tout en
préservant l'intégrité référentielle.

Qdrant indexe les vecteurs avec une structure \textbf{HNSW} (\emph{Hierarchical
Navigable Small World}), qui offre une recherche par plus proches voisins
approchés en temps quasi logarithmique~\cite{malkov2018hnsw}, condition de la
scalabilité de la récupération. La métrique de similarité retenue est le cosinus,
adaptée aux \emph{embeddings} sémantiques normalisés.

\section{Pipeline de récupération et reclassement}
La récupération suit une stratégie en deux temps (figure~\ref{fig:pipe-rag}).
D'abord, une recherche large remonte un ensemble de candidats (paramètre
\code{retrieval\_candidates}, par défaut~15) combinant similarité sémantique et
appariement lexical (récupération \emph{hybride}). Ensuite, un
\textbf{reclassement} (\emph{reranking}) réordonne ces candidats par pertinence
fine et n'en retient que les meilleurs (\code{rerank\_top\_k}, par défaut~5). Ce
schéma « rappel large puis précision fine » maximise la qualité du contexte fourni
au générateur tout en bornant sa taille (\code{rag\_max\_context\_chars}).

\section{Pipeline de génération : le flux agentique}
La génération n'est pas un appel monolithique à un LLM, mais un
\emph{flux agentique} (\code{AgenticRAGWorkflow}) de sept agents coopérants,
chacun persisté comme une exécution inspectable (figure~\ref{fig:pipe-agents}).

\begin{figure}[H]
\centering
\resizebox{\textwidth}{!}{%
\begin{tikzpicture}[node distance=4.5mm,>=Stealth]
  \node[stage] (r)  {Router\\\tiny valide \& route};
  \node[stage,right=of r] (qr) {Query\\Rewriter};
  \node[stage,right=of qr] (re) {Retriever\\\tiny Qdrant};
  \node[stage,right=of re] (rk) {Reranker};
  \node[stage,right=of rk] (g)  {Generator\\\tiny LLM};
  \node[stage,right=of g]  (ev) {Evaluator};
  \node[stage,right=of ev] (co) {Corrector};
  \foreach \a/\b in {r/qr,qr/re,re/rk,rk/g,g/ev,ev/co}{ \draw[flow] (\a) -- (\b); }
  \draw[flow,dashed] (ev.south) to[out=-90,in=-90] node[below,font=\tiny]{si non fidèle} (co.south);
  \node[font=\tiny\itshape,below=10mm of rk,color=slate]
    {Chaque agent $\rightarrow$ une ligne \code{agent\_runs} (statut, latence, entrée/sortie)};
\end{tikzpicture}}
\caption{Pipeline de génération agentique (7 agents)}
\label{fig:pipe-agents}
\end{figure}

\begin{figure}[H]
\centering
\resizebox{0.9\textwidth}{!}{%
\begin{tikzpicture}[node distance=4mm,>=Stealth]
  \node[stage,fill=tierApi] (q) {Question};
  \node[stage,fill=tierAI,right=12mm of q] (cand) {Candidats\\(k=15, hybride)};
  \node[stage,fill=tierAI,right=12mm of cand] (top) {Top 5\\reclassés};
  \node[stage,fill=tierClient,right=12mm of top] (ans) {Réponse\\ancrée + citée};
  \draw[flow] (q) -- node[above,font=\tiny]{récupération} (cand);
  \draw[flow] (cand) -- node[above,font=\tiny]{reranking} (top);
  \draw[flow] (top) -- node[above,font=\tiny]{génération} (ans);
\end{tikzpicture}}
\caption{Stratégie « rappel large puis précision fine » de la récupération}
\label{fig:pipe-rag}
\end{figure}

\section{Pipeline d'évaluation et cycle de vie de la qualité}
Après génération, l'agent évaluateur note la réponse selon plusieurs dimensions :
\textbf{fidélité} (la réponse est-elle étayée par le contexte ?),
\textbf{pertinence} (répond-elle à la question ?), \textbf{précision du contexte}
et \textbf{score d'hallucination}. La méthode d'évaluation est explicitement
tracée (\code{llm}, \code{heuristic} ou \code{fallback}), de sorte que la
confiance dans la mesure est elle-même observable. Lorsque la fidélité est
insuffisante, l'agent \textbf{correcteur} intervient pour réparer la réponse, et
la différence entre version générée et version finale est conservée — c'est le
fondement du \emph{diff} de correction de l'inspecteur de traces.

\section{Architecture de streaming}
Le \emph{streaming} repose sur le protocole \textbf{Server-Sent Events} (SSE) :
le \textit{backend} renvoie une \code{StreamingResponse} FastAPI émettant une
séquence d'événements typés (\code{trace}, \code{citations}, \code{token},
\code{evaluation}, \code{message}, \code{done}). Côté \textit{frontend}, le flux
est consommé via l'API \code{fetch} et un \code{ReadableStream} décodé
incrémentalement. Ce choix — plutôt que des WebSockets — privilégie la simplicité
d'un flux unidirectionnel serveur\,$\rightarrow$\,client, suffisant pour la
diffusion de réponses et naturellement compatible avec HTTP/2.

\section{Architecture d'observabilité : RAGOps et traçabilité}
L'observabilité est traitée comme une fonctionnalité de premier ordre, à trois
niveaux complémentaires :
\begin{enumerate}[leftmargin=*]
  \item \textbf{Métriques système} — Prometheus collecte les métriques du
  \textit{backend} FastAPI (support multi-processus), restituées dans des tableaux
  de bord Grafana provisionnés.
  \item \textbf{RAGOps} — une tour de contrôle agrège la santé des pipelines :
  état de la collection Qdrant, file de traitement documentaire, documents en
  échec et possibilité de relance.
  \item \textbf{Traçabilité applicative} — un explorateur liste les traces de
  requêtes, et un \textbf{inspecteur} reconstruit le cycle de vie complet d'une
  requête (chronologie des agents, passages récupérés, décisions de reclassement,
  citations, évaluation et \emph{diff} du correcteur).
\end{enumerate}

\section{Architecture de sécurité}
La sécurité est multicouche. L'\textbf{authentification} repose sur des jetons
JWT signés ; l'\textbf{autorisation} s'appuie sur un contrôle d'accès par rôles
(utilisateur, admin, super-admin) appliqué au niveau des dépendances FastAPI. Les
données sont \textbf{isolées par espace de travail} : toute requête est cantonnée
au périmètre des espaces dont l'utilisateur est membre. Les actions
d'administration sont \textbf{journalisées} (table \code{audit\_logs}). Enfin, les
vues de débogage \textbf{masquent systématiquement les secrets} (clés contenant
\emph{password}, \emph{secret}, \emph{token}, \emph{api\_key}), garantissant
qu'aucune information sensible ne fuit par l'observabilité.

\section{Architecture de scalabilité, résilience et performance}
\paragraph{Scalabilité.} Le \textit{backend} FastAPI est \emph{sans état} :
plusieurs instances peuvent être placées derrière un répartiteur de charge. Les
\emph{workers} Celery se multiplient horizontalement pour absorber les pics
d'ingestion. Qdrant supporte le partitionnement de collection. PostgreSQL peut
être répliqué en lecture.

\paragraph{Stratégie de \emph{cache}.} Redis joue un double rôle : courtier de
messages pour Celery et magasin de résultats. Il constitue également le socle
naturel d'une mise en cache applicative (réponses, \emph{embeddings} récurrents).

\paragraph{Résilience et gestion des défaillances.} Trois mécanismes assurent la
robustesse : la \textbf{relance} automatique des tâches d'ingestion (jusqu'à trois
tentatives) ; le \textbf{repli} de modèle de génération (\emph{fallback} de Gemini
2.5 Flash vers une version antérieure) ; et l'\textbf{agent correcteur} qui répare
les réponses non fidèles plutôt que de les laisser sortir telles quelles.

\paragraph{Optimisation des performances.} L'asynchronisme bout-en-bout (E/S non
bloquantes côté API, déport du calcul lourd côté Celery), la recherche HNSW
sous-linéaire et le \emph{streaming} (réduisant le \emph{time-to-first-token})
concourent à une expérience réactive.

\section{Justifications technologiques et compromis}
Chaque brique a été retenue à l'issue d'une analyse de compromis explicite,
synthétisée dans le tableau~\ref{tab:choix}.

\begin{table}[H]
\centering
\caption{Décisions technologiques et justifications}
\label{tab:choix}
\small
\begin{tabularx}{\textwidth}{@{}>{\bfseries}l X@{}}
\toprule
Choix & Justification et compromis \\
\midrule
FastAPI &
Framework asynchrone hautes performances, validation Pydantic native, support
direct du \emph{streaming} (SSE). Compromis : écosystème plus jeune que des
frameworks synchrones établis. \\
PostgreSQL &
SGBD relationnel robuste, transactionnel, idéal pour l'état et les traces
structurées ; support JSON pour les champs semi-structurés. Compromis : moins
adapté au stockage vectoriel natif d'où le recours à Qdrant. \\
Qdrant &
Base vectorielle spécialisée, index HNSW performant, API simple, conteneurisable
et auto-hébergeable. Compromis : composant additionnel à exploiter. \\
Redis &
Courtier de messages à faible latence et magasin de résultats pour Celery, base
de \emph{cache}. Compromis : persistance mémoire à dimensionner. \\
Celery &
Standard éprouvé du traitement de tâches asynchrones en Python, relances et
montée en charge des \emph{workers}. Compromis : complexité opérationnelle accrue. \\
Next.js &
\emph{App Router} React 19, rendu hybride, excellente expérience développeur,
\emph{streaming} natif côté client. Compromis : couplage fort à l'écosystème React. \\
Ollama / Gemini &
Abstraction multi-fournisseurs : Ollama pour des \emph{embeddings} locaux et un
fonctionnement hors-ligne, Gemini pour une génération de haute qualité.
Compromis : variabilité de comportement entre fournisseurs, gérée par l'agent
évaluateur. \\
\bottomrule
\end{tabularx}
\end{table}

\section{Conclusion du chapitre}
L'architecture de \devpilot{} traduit ses principes directeurs en un système
concret, en couches, asynchrone, traçable et neutre vis-à-vis des fournisseurs
d'IA. Les pipelines documentaire, de récupération, de génération agentique et
d'évaluation forment un tout cohérent, instrumenté de bout en bout. Le chapitre
suivant détaille la réalisation de ces composants.

%==============================================================================
\chapter{Technologies et réalisation}
\label{chap:real}
%==============================================================================

\section{Présentation des technologies}
La réalisation s'appuie sur une pile technologique moderne et cohérente,
synthétisée dans le tableau~\ref{tab:stack}. L'ensemble est orchestré par Docker
Compose, garantissant la reproductibilité de l'environnement.

\begin{table}[H]
\centering
\caption{Pile technologique de DevPilot AI}
\label{tab:stack}
\small
\begin{tabularx}{\textwidth}{@{}>{\bfseries}l l X@{}}
\toprule
Niveau & Technologie & Rôle \\
\midrule
Frontend     & Next.js, React, TypeScript, Tailwind & Interface, chat, tableaux de bord \\
             & React Query, Recharts & État serveur et visualisations \\
Backend      & FastAPI, Pydantic & API REST et \emph{streaming}, validation \\
             & SQLAlchemy, Alembic & ORM et migrations de schéma \\
Base de données & PostgreSQL 16 & État relationnel et traces \\
Messagerie   & Redis 7 & Courtier de messages et \emph{cache} \\
Asynchrone   & Celery & Traitement documentaire en arrière-plan \\
Vectoriel    & Qdrant 1.12 & Recherche par plus proches voisins (HNSW) \\
Modèles IA   & Ollama, Gemini & \emph{Embeddings} locaux, génération \\
Observabilité& Prometheus, Grafana & Métriques et tableaux de bord \\
Déploiement  & Docker, Docker Compose & Conteneurisation et orchestration \\
\bottomrule
\end{tabularx}
\end{table}

\section{Réalisation du \textit{frontend}}
L'interface est bâtie avec Next.js (\emph{App Router}) et un système de design
maison fondé sur des jetons CSS (thème clair/sombre), garantissant cohérence et
accessibilité. L'état serveur est géré par React Query, qui mutualise le cache et
la synchronisation des requêtes. Le listing~\ref{lst:rq} illustre le \emph{hook}
d'inspection de trace, qui orchestre quatre requêtes parallèles.

\begin{lstlisting}[language=JavaScript,caption={Hook React Query d'inspection de trace (extrait)},label=lst:rq]
export function useTraceInspector(messageId, includeContent, enabled) {
  const trace = useQuery({
    queryKey: ["trace", messageId, includeContent],
    queryFn: () => getRagTrace(messageId, includeContent),
    enabled,
  });
  // Requetes secondaires en parallele, degradation gracieuse
  const retrieval  = useQuery({ queryKey: ["trace-retrieval", messageId], ... });
  const reranking  = useQuery({ queryKey: ["trace-reranking", messageId], ... });
  const evaluation = useQuery({ queryKey: ["trace-evaluation", messageId], ... });
  return { trace, retrieval, reranking, evaluation, /* ... */ };
}
\end{lstlisting}

\figplaceholder{Figure 4.1 — Interface du chat conversationnel avec réponse en
\emph{streaming}, citations en ligne et cartes de sources.}

\section{Réalisation du \textit{backend}}
Le \textit{backend} FastAPI expose une API versionnée organisée en routeurs
(authentification, espaces, documents, chat, retrieval, RAGOps, traces, santé).
La logique métier est encapsulée dans des services dédiés
(\code{document\_service}, \code{retrieval\_service}, \code{reranking\_service},
\code{evaluation\_service}, \code{rag\_traces\_service}, \code{ragops\_service}),
séparant clairement les routes des règles du domaine.

\figplaceholder{Figure 4.2 — Documentation interactive de l'API (Swagger/OpenAPI)
générée automatiquement par FastAPI.}

\section{Réalisation de la base de données}
Le schéma relationnel (onze tables) est défini avec SQLAlchemy et son évolution
est gérée par des migrations Alembic versionnées (schéma initial, évaluation des
réponses, RBAC d'administration). Les identifiants sont des UUID, et des
\emph{mixins} factorisent les colonnes temporelles communes.

\section{Traitement documentaire}
La tâche \code{process\_document\_task} (Celery) orchestre l'extraction du texte,
le découpage, la vectorisation et l'indexation Qdrant, avec relance automatique en
cas d'échec et passage du document au statut \code{failed} le cas échéant — état
exploité par RAGOps pour proposer une relance.

\figplaceholder{Figure 4.3 — Interface de gestion documentaire : glisser-déposer,
téléversement par lots, indicateurs de progression et de statut.}

\section{Pipeline RAG}
Le pipeline agentique (\code{AgenticRAGWorkflow}) instancie les sept agents et
journalise chaque exécution dans \code{agent\_runs}. Cette persistance
systématique est ce qui rend l'inspecteur de traces possible : aucun état n'est
implicite.

\section{Chat en \emph{streaming}}
La diffusion en \emph{streaming} relie une \code{StreamingResponse} SSE côté
serveur à un consommateur \code{ReadableStream} côté client. Les événements
\code{trace} et \code{citations} arrivent avant les jetons de texte, permettant à
l'interface d'afficher les sources dès le début de la génération.

\figplaceholder{Figure 4.4 — Réponse en cours de \emph{streaming} avec indicateur
de saisie, défilement automatique et suggestions de questions de suivi.}

\section{Tableau de bord RAGOps}
La tour de contrôle RAGOps agrège la santé de la collection Qdrant, l'état de la
file d'ingestion, les documents en échec et les indicateurs de qualité agrégés.
Elle constitue le point d'entrée de l'exploitation quotidienne de la plateforme.

\figplaceholder{Figure 4.5 — Tour de contrôle RAGOps : santé des pipelines, file
de traitement et alertes.}

\section{Explorateur et inspecteur de traces}
L'explorateur liste les traces avec filtres serveur (recherche, espace, stratégie
de récupération, seuils de fidélité et d'hallucination, plage de dates).
L'inspecteur reconstruit, pour une trace donnée, neuf sections : résumé,
chronologie des agents, explorateur de récupération, explorateur de reclassement,
explorateur de citations, centre d'évaluation, analyse du correcteur (avec
\emph{diff} mot à mot), usage des modèles et graphe interactif du pipeline. Un
score de santé de trace (sur~100) synthétise la qualité d'une requête.

\figplaceholder{Figure 4.6 — Inspecteur de traces : chronologie des agents,
passages récupérés, citations et scores d'évaluation.}

\section{Interface d'administration}
Les pages d'administration (utilisateurs, espaces, documents, journaux d'audit,
paramètres) adoptent une ergonomie inspirée des consoles d'exploitation
professionnelles : métriques visuelles, indicateurs de santé, recherche, filtres
et actions groupées, plutôt qu'un CRUD austère.

\figplaceholder{Figure 4.7 — Console d'administration : métriques, filtres et
actions groupées sur les entités de la plateforme.}

\section{Sécurité}
La réalisation applique les principes du chapitre~3 : génération et vérification
de jetons JWT, dépendances d'autorisation par rôle, isolation des données par
espace, journalisation d'audit et masquage des secrets dans toute vue de
débogage.

\section{Tests et performance}
La réalisation s'accompagne de tests automatisés (unitaires, intégration, API) et
d'une attention portée à la performance, détaillés au chapitre suivant.

\section{Conclusion du chapitre}
La réalisation concrétise fidèlement l'architecture conçue : un \textit{frontend}
réactif et typé, un \textit{backend} en couches asynchrone, une ingestion
documentaire robuste, un pipeline agentique traçable et un sous-système
d'observabilité complet.

%==============================================================================
\chapter{Tests et validation}
\label{chap:tests}
%==============================================================================

\section{Stratégie de test}
La validation suit une pyramide de tests classique, des tests unitaires rapides
et nombreux jusqu'aux tests d'acceptation de plus haut niveau. Le
tableau~\ref{tab:tests} synthétise les catégories couvertes.

\begin{table}[H]
\centering
\caption{Catégories de tests et objets validés}
\label{tab:tests}
\small
\begin{tabularx}{\textwidth}{@{}>{\bfseries}l X@{}}
\toprule
Catégorie & Objet de la validation \\
\midrule
Tests unitaires & Fonctions pures : découpage, scores, \emph{diff} mot à mot, extraction d'usage \\
Tests d'intégration & Chaîne ingestion $\rightarrow$ indexation $\rightarrow$ récupération \\
Tests d'API & Contrats des \emph{endpoints} (statuts, schémas, autorisation) \\
Validation RAG & Pertinence des passages, ancrage et présence des citations \\
Tests de performance & Latence des requêtes, débit d'ingestion, recherche vectorielle \\
Tests de \emph{streaming} & Ordre et complétude des événements SSE \\
Tests de sécurité & Isolation par espace, contrôle d'accès par rôles, masquage des secrets \\
\bottomrule
\end{tabularx}
\end{table}

\section{Tests unitaires}
Les fonctions déterministes de la couche de présentation et du domaine sont
couvertes par des tests unitaires : par exemple l'algorithme de \emph{diff} mot à
mot (plus longue sous-séquence commune) utilisé par l'analyse du correcteur,
l'extraction d'usage des modèles, ou le calcul du score de santé de trace.

\section{Tests d'intégration et d'API}
Les tests d'intégration valident la chaîne complète d'ingestion et de
récupération sur un corpus contrôlé, tandis que les tests d'API vérifient les
contrats des \emph{endpoints}, y compris les codes de statut et l'application des
règles d'autorisation.

\section{Validation du RAG}
La qualité du RAG est validée selon trois axes : la \textbf{récupération}
(les passages remontés sont-ils pertinents ?), l'\textbf{ancrage} (la réponse
est-elle étayée par ces passages ?) et la \textbf{citation} (les sources sont-elles
correctement référencées ?). L'agent évaluateur fournit une mesure continue de ces
axes, exposée dans l'inspecteur de traces.

\section{Tests de performance et de \emph{streaming}}
Les tests de performance mesurent la latence de bout en bout, le débit
d'ingestion et la latence de la recherche vectorielle. Les tests de
\emph{streaming} vérifient que les événements SSE arrivent dans l'ordre attendu
(\code{trace}, \code{citations}, \code{token}\dots, \code{done}) et que le flux se
termine proprement.

\section{Tests de sécurité}
Les tests de sécurité valident l'isolation des données entre espaces de travail,
le respect du contrôle d'accès par rôles et l'absence de fuite de secrets dans les
vues de débogage et d'audit.

\section{Critères d'acceptation}
Un scénario est considéré comme validé lorsque : une question sur un corpus indexé
produit une réponse ancrée et citée ; l'évaluation associée est calculée et
persistée ; la trace complète est reconstructible dans l'inspecteur ; et aucune
fuite de données inter-espaces ni de secret n'est observée.

\section{Conclusion du chapitre}
La stratégie de validation couvre l'ensemble des couches, du test unitaire à
l'acceptation, en accordant une attention particulière aux propriétés
spécifiques d'une plateforme d'IA : ancrage, citations, qualité mesurée et
sécurité.

%==============================================================================
\chapter*{Limites et perspectives}
\addcontentsline{toc}{chapter}{Limites et perspectives}
\markboth{Limites et perspectives}{Limites et perspectives}
%==============================================================================

\section*{Limites actuelles}
Le système, bien que de qualité production, présente des limites assumées. La
comptabilité fine d'usage des modèles (jetons, coût) dépend des métadonnées
exposées par chaque fournisseur et n'est pas toujours disponible — l'inspecteur
affiche alors honnêtement « non capturé » plutôt qu'une valeur fabriquée. Le
déploiement cible est mono-machine via Docker Compose, sans orchestration
multi-n{\oe}uds. L'évaluation repose sur des heuristiques et un juge LLM plutôt
que sur un cadre d'évaluation formel à large échelle. Enfin, le multi-tenant reste
au niveau de l'espace de travail, sans cloisonnement organisationnel complet.

\section*{Perspectives}
Plusieurs axes d'évolution se dégagent naturellement :
\begin{itemize}[leftmargin=*]
  \item \textbf{Orchestration agentique avancée} : adoption de LangGraph pour
  modéliser explicitement les transitions et les boucles du flux agentique.
  \item \textbf{Évaluation formelle} : intégration de RAGAS et de jeux
  d'évaluation versionnés pour un suivi longitudinal de la qualité.
  \item \textbf{Déploiement cloud} : migration vers une orchestration Kubernetes
  pour la haute disponibilité et la montée en charge automatique.
  \item \textbf{SaaS multi-tenant} : cloisonnement organisationnel, quotas et
  facturation par tenant.
  \item \textbf{RBAC d'entreprise} : modèle de permissions fin, intégration SSO
  et fédération d'identité.
\end{itemize}

%==============================================================================
\chapter*{Conclusion générale}
\addcontentsline{toc}{chapter}{Conclusion générale}
\markboth{Conclusion générale}{Conclusion générale}
%==============================================================================

Ce projet de fin d'études a abouti à la conception et à la réalisation de
\devpilot{}, une plateforme d'IA d'entreprise pour la gestion des connaissances
d'ingénierie et le question-réponse fondé sur le RAG. L'objectif initial — combler
l'écart entre la puissance générative des LLM et l'exigence de réponses
\emph{fondées}, \emph{citées}, \emph{mesurées} et \emph{traçables} — est atteint.

\paragraph{Objectifs atteints.} La plateforme centralise les connaissances en
espaces isolés, les ingère de façon asynchrone, les indexe vectoriellement, et
restitue des réponses ancrées et citées via un pipeline agentique de sept agents,
le tout diffusé en \emph{streaming}. La qualité de chaque réponse est évaluée et,
si nécessaire, corrigée. L'ensemble est instrumenté par un sous-système
d'observabilité — RAGOps, explorateur et inspecteur de traces — qui distingue
\devpilot{} d'un simple assistant conversationnel.

\paragraph{Contributions techniques.} Les contributions d'ingénierie majeures
sont : l'élévation de la \emph{trace} au rang de concept de premier ordre du
modèle de données ; un pipeline agentique entièrement persisté et inspectable ;
une architecture asynchrone, en couches et neutre vis-à-vis des fournisseurs d'IA ;
et un inspecteur de traces reconstruisant le cycle de vie complet d'une requête.

\paragraph{Enseignements d'ingénierie.} Ce travail a confirmé que la valeur d'une
plateforme d'IA d'entreprise ne réside pas seulement dans la qualité des réponses,
mais dans leur \emph{auditabilité} et leur \emph{exploitabilité}. Concevoir pour
l'observabilité dès l'origine, traiter les exigences non fonctionnelles comme des
exigences de premier ordre, et raisonner systématiquement en termes de compromis
ont été les leviers déterminants de ce projet.

\paragraph{Impact professionnel.} \devpilot{} démontre une maîtrise de
l'architecture des systèmes distribués, de l'ingénierie des plateformes d'IA et
des pratiques d'exploitation modernes, au niveau attendu d'un ingénieur
concepteur de systèmes plutôt que d'un simple prototype académique.

%==============================================================================
\begin{thebibliography}{99}
\addcontentsline{toc}{chapter}{Bibliographie}

\bibitem{lewis2020rag} P.~Lewis \emph{et al.}, « Retrieval-Augmented Generation
for Knowledge-Intensive NLP Tasks », \emph{NeurIPS}, 2020.

\bibitem{karpukhin2020dpr} V.~Karpukhin \emph{et al.}, « Dense Passage Retrieval
for Open-Domain Question Answering », \emph{EMNLP}, 2020.

\bibitem{robertson2009bm25} S.~Robertson et H.~Zaragoza, « The Probabilistic
Relevance Framework: BM25 and Beyond », \emph{Foundations and Trends in
Information Retrieval}, 2009.

\bibitem{malkov2018hnsw} Y.~Malkov et D.~Yashunin, « Efficient and Robust
Approximate Nearest Neighbor Search Using Hierarchical Navigable Small World
Graphs », \emph{IEEE TPAMI}, 2018.

\bibitem{reimers2019sbert} N.~Reimers et I.~Gurevych, « Sentence-BERT: Sentence
Embeddings Using Siamese BERT-Networks », \emph{EMNLP}, 2019.

\bibitem{nogueira2019rerank} R.~Nogueira et K.~Cho, « Passage Re-ranking with
BERT », \emph{arXiv:1901.04085}, 2019.

\bibitem{es2023ragas} S.~Es \emph{et al.}, « RAGAS: Automated Evaluation of
Retrieval Augmented Generation », \emph{arXiv:2309.15217}, 2023.

\bibitem{fastapi} S.~Ramírez, \emph{FastAPI Documentation},
\url{https://fastapi.tiangolo.com}, consulté en 2026.

\bibitem{qdrant} \emph{Qdrant — Vector Database Documentation},
\url{https://qdrant.tech/documentation}, consulté en 2026.

\bibitem{celery} \emph{Celery — Distributed Task Queue Documentation},
\url{https://docs.celeryq.dev}, consulté en 2026.

\bibitem{redis} \emph{Redis Documentation}, \url{https://redis.io/docs}, consulté
en 2026.

\bibitem{postgres} \emph{PostgreSQL 16 Documentation},
\url{https://www.postgresql.org/docs/16}, consulté en 2026.

\bibitem{nextjs} \emph{Next.js Documentation}, \url{https://nextjs.org/docs},
consulté en 2026.

\bibitem{docker} \emph{Docker and Docker Compose Documentation},
\url{https://docs.docker.com}, consulté en 2026.

\bibitem{gemini} Google DeepMind, \emph{Gemini API Documentation},
\url{https://ai.google.dev}, consulté en 2026.

\bibitem{ollama} \emph{Ollama Documentation}, \url{https://ollama.com}, consulté
en 2026.

\end{thebibliography}

\end{document}
